\section{讨论与局限}
\label{sec:discussion}

\subsection{78.7\% 对无人值守 CI 自愈意味着什么}
LLM 自愈最常见的部署模式是高度自主化的：一个失败的端到端测试触发 LLM 给出补丁，CI 自动提交补丁或将其作为\rev{拉取请求（Pull Request，PR）}评论发出。由 \S\ref{sec:f1} 的对抗性上界粗略换算，在不带\rev{行为验证器}的此类部署中，约 15\%（粗略上界估计，非实测部署率）的真实 Playwright 断裂存在被朴素修复器静默改写到错误元素的风险。该风险的危险并不在于“修复器让构建失败”——构建失败本就是测试应承担的职责——而在于“修复器使构建通过，但通过对的是错误的目标”。

\subsection{为什么软提示安全网会失效}
\textsc{Vanilla} 与 \textsc{Abstain} 两种提示之间的等同性（表~\ref{tab:falseheal}）并不令人意外。一种可能的解释是：系统提示中“必要时放弃”的软指令需与“生成合理输出”这一更强先验竞争，多数情况下处于劣势（此为机理层面的推测，本文未做训练动力学层面的验证）。对工具设计的启示是清晰的：反误修复的安全保障必须置于模型\emph{之外}，例如确定性的检索分数阈值、锚点类别归属规则，以及一个能够证伪通过补丁的行为验证器。

\subsection{为什么两个低开销确定性杠杆胜过提示工程}
\S\ref{sec:f2} 的 testId 加权检索与强锚点后置过滤在机制上显而易见。经二因子析因分离，二者各带来 +2、合计 +3 个精确匹配且零回归（表~\ref{tab:l3-factorial}）；增益方向一致，但在 $n=75$ 下经 McNemar 精确检验未达统计显著。因此，文章既不声称其统计显著，也不声称该增益能泛化到当前设置以外，而是将其作为 HealReact 的一条局部设计经验：把锚点优先级约束放到候选池里和输出后置过滤里，比依靠系统提示让模型把约束内化更可靠，且不引入回归。

\subsection{局限}
\textit{模型单一性带来的有效性威胁。}本文全部数字基于单一开源小型代码模型（\texttt{qwen2.5-coder:7b}）。前沿大模型（如 GPT-4 类）是否能在隐性误修复维度上显著缩小该差距尚待实测；根据 \textsc{Vanilla = Abstain} 等同性所揭示的机理，文章倾向于认为单凭规模不足以在缺乏外部确定性验证器的前提下闭合该差距。
\textit{基准数据集范围。}主结果来自单一 React+Playwright 代码库 koenig；\S\ref{sec:f3} 中已扩展至 payload，但其可达率受到 BEM 类名盲区的限制。
\textit{静态修复代理指标与运行时通过率。}77.3\% / 62.4\% 这两个 L3 数字反映的是选择器解析器层面的精确一致性。其与真实 Playwright 运行时通过率之间的换算关系，需要在 ReproBreak 提供的逐\rev{版本提交}复现环境中通过 Docker 重放\rev{修复后的测试}加以建立。作为最小可行性检查，文章对 3 个 v1 静态精确匹配用例（\rev{编号 615、702、703}）做了 Docker 运行时重放：三者的\rev{修复基线}均通过；替换为 HealReact 提议选择器后，2/3 运行通过，1/3 运行失败（\rev{编号 702}，等待 \texttt{getByTestId('color-picker-toggle')} 超时）。这一小规模验证实验不能估计全量通过率，但说明静态精确匹配既有运行时可迁移的案例，也会因可见性/时序/上下文差异而失败；因此本文继续把 77.3\% / 62.4\% 标为静态修复代理指标，而不宣称端到端通过率。

\subsection{有效性威胁}
78.7\% 这一数字依赖解析器语义。\texttt{resolve\_locators.py} 把动态 JS 表达式属性值视为可能匹配、对复合 CSS 使用祖先过近似；这两项决策在错误排查场景下合理，但在边界情况下偏松。文章在解析器假设审计中撤回了早期“可达率上限 98\%”的虚抬声明，在收紧解析器假设后以 80.6\% 作为保守数字。由于解析器是本文所有“正确匹配/误修复/可达率”判定的\rev{事实判定器}，其判定边界会直接影响这些数字，文章做了一次 strict/permissive 双假设的敏感性扫描（表~\ref{tab:resolver-sens}）。本文主要实验结果中的 80.6\%（75/93）来自 permissive 假设；若关闭上述两处过近似改用 strict 假设，可达率降至 54.8\%（51/93）。因此 L1 静态可达率的真实区间应理解为 [54.8\%, 80.6\%]，80.6\% 是上界。此外，permissive 下 75 个可达用例中有 46 个为多匹配（解析器按 \texttt{(componentFile, line)} 顺序取首个命中消歧），未解析用例从 permissive 的 18 个增至 strict 的 42 个。对抽样用例用真实 DOM 快照交叉校验解析判定，仍留作后续工作。

\begin{table}[t]
\centering
\caption{解析器 strict/permissive 敏感性（koenig，$n=93$）。permissive 为本文所采用的假设；strict 关闭“动态属性值视为可能匹配”与“同文件祖先近似”两处过近似。}
\label{tab:resolver-sens}
\small
\resizebox{\columnwidth}{!}{%
\begin{tabular}{lrr}
\toprule
指标 & permissive（本文） & strict（下界） \\
\midrule
new\_locator 可达 & 75 / 93 (80.6\%) & 51 / 93 (54.8\%) \\
多匹配（$\ge2$ 命中） & 46 & 0 \\
未解析（0 命中） & 18 & 42 \\
\bottomrule
\end{tabular}}
\end{table}

\subsection{未来工作}
本文的诊断结论自然延伸出两类可推进的工程化方向：其一，将静态基底向跨锚点文化（BEM 模板字面量、CSS Modules 等）扩展，使可达率在不同 React 项目中保持稳定；其二，沿 \S\ref{sec:findings} 给出的经验性必要性，构建并集成一个基于网络行为重放的轻量级验证器，作为 LLM 修复结果的下游守门。
